% =====================================================================
% RAPPORT TECHNIQUE SCIENTIFIQUE --- SENTINELBAND (v1.3)
% Version recentrée sur le prototype wearable (partie antenne retirée)
%
% Images optionnelles (placer dans le même dossier que le .tex) :
%   fig-mpu6050-chute.png  : enregistrement réel annoté (g en fonction du temps)
%   photo-prototype.jpg    : photo du montage
%   capture-app.png        : captures de l'application mobile
%   log-serie.png          : extrait des logs série
% Si un fichier est absent, un cadre rouge « À COMPLÉTER » s'affiche.
% Avant le rendu : supprimer tous les \aCompleter{...} et ce bloc de commentaires.
% =====================================================================
\documentclass[12pt,a4paper]{article}

% --- Encodage et langue ---
\usepackage[utf8]{inputenc}
\usepackage[T1]{fontenc}
\usepackage[french]{babel}
\usepackage{lmodern}
\usepackage{eurosym}

% --- Mise en page ---
\usepackage[top=2.5cm, bottom=2.5cm, left=2.5cm, right=2.5cm]{geometry}
\usepackage{setspace}
\onehalfspacing

% --- Mathématiques ---
\usepackage{amsmath, amssymb}

% --- Tableaux ---
\usepackage{booktabs}
\usepackage{array}
\usepackage{tabularx}
\usepackage{multirow}

% --- Figures ---
\usepackage{graphicx}
\usepackage{float}
\usepackage[font=small,labelfont=bf]{caption}
\usepackage{tikz}
\usetikzlibrary{arrows.meta,positioning}

% --- Couleurs ---
\usepackage[dvipsnames]{xcolor}

% --- Code source ---
\usepackage{listings}
\lstset{
  basicstyle=\ttfamily\footnotesize,
  breaklines=true,
  frame=single,
  language=C++,
  commentstyle=\color{gray},
  keywordstyle=\color{blue},
  stringstyle=\color{red!70!black}
}

% --- Bibliographie ---
\usepackage[numbers,sort&compress]{natbib}

% --- En-têtes ---
\usepackage{fancyhdr}
\pagestyle{fancy}
\fancyhf{}
\fancyhead[L]{\small Rapport technique \textemdash{} SentinelBand}
\fancyhead[R]{\small\thepage}
\fancyfoot[C]{\small Année universitaire 2026\textendash 2027}

% --- Liens (en dernier) ---
\usepackage[colorlinks=true, linkcolor=blue!70!black, citecolor=green!50!black, urlcolor=blue!60!black]{hyperref}

% --- Commandes personnalisées ---
\newcommand{\aCompleter}[1]{\textcolor{red}{\textbf{[À COMPLÉTER : #1]}}}
\newcommand{\aRemplir}{\textcolor{red}{\textbf{[n]}}}

% Autorise une coupure de ligne après chaque underscore (identifiants longs en \texttt)
\renewcommand{\_}{\textunderscore\allowbreak}
\emergencystretch=3em

% Figure optionnelle : #1 fichier, #2 largeur (fraction de \textwidth), #3 légende, #4 label
\newcommand{\figureopt}[4]{%
\begin{figure}[H]
  \centering
  \IfFileExists{#1}%
    {\includegraphics[width=#2\textwidth]{#1}}%
    {\fbox{\parbox{0.8\textwidth}{\centering\vspace{1cm}\aCompleter{insérer le fichier #1}\vspace{1cm}}}}
  \caption{#3}
  \label{#4}
\end{figure}}

% =====================================================================
%                         PAGE DE GARDE
% =====================================================================
% =====================================================================
%                         PAGE DE GARDE
% =====================================================================
\begin{document}

\begin{titlepage}
  \centering

\thispagestyle{empty}

{\Large \textbf{ESPRIT -- École Supérieure Privée d’Ingénierie et de Technologies}\par}
  \vspace{0.4cm}
\begin{figure}
    \centering
    \includegraphics[width=0.5\linewidth]{Logo_ESPRIT_Ariana.jpg}
   
    \label{fig:placeholder}
\end{figure}
\centering



\vspace{1cm}

\rule{\linewidth}{0.4pt}\par
\vspace{0.2cm}
{\Large \textbf{Conception et évaluation préliminaire d'un ceinture
connecté pour la surveillance de la sécurité des travailleurs industriels~: prototype wearable multi-capteurs avec algorithme de détection de chute}\par}
\vspace{0.2cm}
\rule{\linewidth}{0.4pt}\par

\vspace{0.6cm}


\vspace{0.8cm}

\small
\begin{tabular}{r l}
  \textbf{Étudiants}   & Adem Aboueb \\
                        & Aziz Ayari \\
                        & Nesrine Fezzani \\
                        & Mohamed Khalil Hammami \\
                        & Mohamed Wassef Mnasri \\[4pt]
  \textbf{Encadrante}   & Ameni Mersani, Wearable Technologies, ESPRIT \\[4pt]
  \textbf{Formation}    & 5\textsuperscript{ème} année SLEAM --- Wearable Technologies \\[4pt]
  \textbf{Version}      & v1.3 --- 22 octobre 2026 \\[4pt]
  \textbf{Dépôt du code} & \url{https://github.com/AdemAboueb7/sentinelband.git}
\end{tabular}


\vfill
{\small Année universitaire 2026--2027\par}
\end{titlepage}

% =====================================================================
%                            RÉSUMÉ
% =====================================================================
\newpage
\section*{Résumé}
\addcontentsline{toc}{section}{Résumé}

La sécurité des travailleurs isolés en milieu industriel constitue un enjeu majeur, les chutes pouvant entraîner des conséquences graves en l'absence d'alerte rapide. Ce travail présente SentinelBand, un prototype de ceinture connecté fondé sur un ESP32 intégrant un accéléromètre-gyroscope MPU6050 ($\pm$8\,g, $\pm$250\,°/s), un récepteur GPS NEO-6M, un module GSM SIM800L, un buzzer, une LED RGB et un bouton d'annulation. Un algorithme embarqué de détection de chute à machine à états finis (FSM) exploite quatre phases (pré-chute, apesanteur, impact, posture post-chute) avec un score de confiance pondéré (0--100). Un score $\geq$60 déclenche des alertes graduées (jaune puis rouge), transmises par Wi-Fi à une API REST Node.js/PostgreSQL et par SMS au responsable de sécurité. Une application mobile React Native assure la supervision en temps réel, avec notifications push, historique et géolocalisation. Les essais par injection de chute via la commande série ont validé la chaîne d'alerte et la transmission vers l'API~; l'algorithme de détection lui-même n'a pas encore été validé sur des chutes réelles. Les principales limites sont l'absence de capteurs de température et de gaz réels, l'absence de validation statistique, l'échantillonnage à 50\,Hz, l'absence de mesure d'autonomie et la transmission non chiffrée.



% =====================================================================
%                     TABLE DES MATIÈRES
% =====================================================================
\newpage
\tableofcontents
\newpage

% =====================================================================
%              I. INTRODUCTION
% =====================================================================
\section{Introduction}

\subsection{Contexte des systèmes wearables}

Les dispositifs électroniques portables (\textit{wearables}) connaissent un essor considérable dans les domaines de la santé, du sport et de la sécurité au travail~\cite{ref_wearable_survey}. Dans le secteur industriel, la surveillance des travailleurs isolés ou exposés à des environnements dangereux (chantiers, sites pétrochimiques, entrepôts logistiques) représente un besoin critique~\cite{ref_lone_worker}. Les systèmes de type \textit{man-down} ou \textit{lone worker protection} combinent typiquement des capteurs inertiels pour la détection de chute, des capteurs environnementaux (température, gaz), un module de géolocalisation et un canal de communication sans fil pour transmettre les alertes en temps réel. La miniaturisation des microcontrôleurs connectés tels que l'ESP32~\cite{ref_esp32_datasheet}, couplée à la disponibilité de capteurs MEMS bon marché comme le MPU6050~\cite{ref_mpu6050_datasheet}, rend accessibles des prototypes wearables complets pour un coût de composants de l'ordre de quelques dizaines d'euros.

\subsection{Problème scientifique}

Comment concevoir un système embarqué portable multi-capteurs capable de détecter de manière fiable les situations de danger (chute, immobilité prolongée) d'un travailleur industriel et de transmettre des alertes en temps réel vers une application mobile de supervision, tout en limitant les faux positifs susceptibles de perturber le travail~?

\subsection{État de l'art}

La détection de chute par accéléromètre embarqué fait l'objet d'une littérature abondante. Bayar \textit{et al.}~\cite{ref_bayar2023} ont déployé un modèle de classification de chute sur un ESP32 couplé à un MPU6050 via la plateforme Edge Impulse, illustrant l'approche TinyML sur ce type de matériel. Les approches à seuils exploitent la magnitude du vecteur accélération $|\vec{a}| = \sqrt{a_x^2 + a_y^2 + a_z^2}$, éventuellement combinée à la vitesse angulaire et à la posture après impact~\cite{ref_fall_detection_threshold}. Kangas \textit{et al.}~\cite{ref_kangas} ont comparé plusieurs algorithmes de faible complexité fondés sur des accéléromètres portés sur le corps et montré l'intérêt de combiner des critères de chute libre, d'impact et de posture finale. Mauldin \textit{et al.}~\cite{ref_mauldin} proposent, à l'inverse, une approche par apprentissage profond sur smartwatch, plus coûteuse en calcul et en données d'entraînement que les méthodes à seuils retenues ici pour un microcontrôleur.

\subsection{Contribution du travail}

Les principales contributions de ce travail sont les suivantes~:
\begin{enumerate}
    \item la conception, l'implémentation et l'évaluation préliminaire d'un prototype wearable SentinelBand intégrant un algorithme de détection de chute à machine à états finis (FSM) en quatre phases avec score de confiance pondéré ($C_{\text{final}} = 0{,}35 \cdot S_{\text{FF}} + 0{,}45 \cdot S_{\text{impact}} + 0{,}20 \cdot S_{\text{posture}}$), sur ESP32 avec MPU6050 ($\pm$8\,g), GPS NEO-6M et module GSM SIM800L~;
    \item le développement d'une chaîne complète de supervision comprenant une API REST (Node.js~18 / PostgreSQL~16) et une application mobile React Native (Expo~57) avec notifications push, gestion de collègues avec consentement explicite et visualisation temps réel de la télémétrie géolocalisée.
\end{enumerate}

% =====================================================================
%         II. SPÉCIFICATIONS ET ARCHITECTURE
% =====================================================================
\section{Spécifications et architecture}

\subsection{Scénario d'usage}

Un travailleur opérant seul sur un chantier industriel porte le dispositif SentinelBand à la ceinture. Le boîtier surveille en continu ses mouvements via le MPU6050 à 50\,Hz. Si une séquence de chute est détectée (apesanteur, impact, posture anormale) avec un score de confiance $\geq 60$, ou si une immobilité prolongée est constatée ($>$20\,s sans mouvement), le système entre en état d'alerte JAUNE~: le buzzer émet un bip intermittent (150\,ms toutes les secondes) et la LED passe au jaune. Si le travailleur n'appuie pas sur le bouton d'annulation dans les 15\,s (\texttt{YELLOW\_TIMEOUT\_MS}), l'état escalade en alerte ROUGE~: un SMS d'urgence contenant la position GPS est envoyé au responsable via le SIM800L, et la télémétrie est transmise à l'API par Wi-Fi toutes les 2\,s. L'application mobile du responsable reçoit une notification push et affiche l'état, la position et les données capteurs. Le travailleur peut annuler une alerte à tout moment (GPIO\,33)~; un délai de refroidissement de 15\,s (\texttt{CANCEL\_HOLDOFF\_MS}), distinct du délai d'escalade, empêche un re-déclenchement immédiat après annulation.

\subsection{Tableau d'exigences vérifiables}



\begin{table}[H]
\centering
\caption{Exigences fonctionnelles et techniques du système SentinelBand.}
\label{tab:exigences}
\footnotesize
\begin{tabularx}{\textwidth}{>{\raggedright\arraybackslash}X >{\raggedright\arraybackslash}X >{\raggedright\arraybackslash}X >{\raggedright\arraybackslash}X >{\raggedright\arraybackslash}X}
\toprule
\textbf{Exigence} & \textbf{Cible} & \textbf{Justification} & \textbf{Vérification} & \textbf{Résultat} \\
\midrule
Acquisition IMU & 50\,Hz, $\pm$8\,g, $\pm$250\,°/s & Détection de chute & Lecture registres I2C 0x3B--0x48 & Testé \\
\addlinespace
Délai d'alerte locale (JAUNE) après le début de la chute & $\leq$8\,s & Réactivité (la fenêtre de posture de 5\,s impose un minimum d'environ 5\,s après l'impact) & Chronométrage série & $\approx$5--7,5\,s (calculé d'après les paramètres FSM, non chronométré) \\
\addlinespace
Délai d'alerte immobilité & $\leq$20\,s & Sécurité & Chronométrage & 20\,s (paramètre, logs) \\
\addlinespace
Connectivité Wi-Fi & IEEE 802.11 b/g/n & Transmission télémétrie & Test HTTP POST & Testé (reconnexion auto) \\
\addlinespace
Alerte hors Wi-Fi & SMS via réseau GSM & Alerte sans infrastructure locale & AT+CMGS & Testé (diagnostic AT, SMS reçu) \\
\addlinespace
Autonomie & $\geq$8\,h & Journée de travail & Estimation & Estimée 8--14\,h, non mesurée \\
\addlinespace
Télémétrie temps réel & 4\,s normal / 2\,s alerte & Supervision continue & Inspection code & Implémenté \\
\addlinespace
Sécurité API & SHA-256 + JWT 30\,j & Confidentialité & Inspection code serveur & Implémenté \\
\bottomrule
\end{tabularx}
\end{table}

\subsection{Architecture fonctionnelle}

Le système se décompose en cinq blocs~:

\begin{enumerate}
    \item \textbf{Bloc capteurs~:} MPU6050 ($\pm$8\,g, $\pm$250\,°/s, I2C à 100\,kHz, adresse 0x68), GPS NEO-6M (UART2, 9600\,bauds, tampon RX 1024 octets). Température et gaz sont des constantes logicielles ($T = 25{,}0$\,°C, gaz~$= 150$\,mV)~: aucun capteur LM35 ni MQ-2 n'est connecté.

    \item \textbf{Bloc traitement/décision~:} ESP32 exécutant la boucle principale (lecture IMU à 50\,Hz, FSM de détection de chute en 4 phases, évaluation d'état NORMAL $\rightarrow$ JAUNE $\rightarrow$ ROUGE, bouton d'annulation, commandes série de débogage).

    \item \textbf{Bloc actionneurs/alerte~:} buzzer piézo (GPIO\,4), LED RGB à cathode commune (GPIO\,25/26/27)~: vert en NORMAL, jaune en JAUNE, rouge en ROUGE~; bouton poussoir (GPIO\,33, \texttt{INPUT\_PULLUP}, debounce 300\,ms).

    \item \textbf{Bloc transmission~:} Wi-Fi intégré de l'ESP32 (antenne PCB interne, mode STA, reconnexion toutes les 8\,s) vers l'API REST (POST JSON, timeouts 2500\,ms et 3000\,ms)~; module SIM800L (UART1) pour le SMS d'urgence géolocalisé en alerte ROUGE.

    \item \textbf{Bloc application~:} API Node.js/Express (port 3000) avec PostgreSQL~16 (Docker Compose), application React Native/Expo avec notifications push (Expo Push Service), polling toutes les 3\,s, historique télémétrique et carte.
\end{enumerate}

\begin{figure}[H]
\centering
\begin{tikzpicture}[
  box/.style={draw, rounded corners, align=center, text width=3.0cm, minimum height=1.1cm, font=\footnotesize},
  esp/.style={draw, rounded corners, align=center, text width=3.0cm, minimum height=5.2cm, font=\footnotesize, fill=blue!8},
  arrow/.style={-{Stealth[length=2.5mm]}, thick}
]
  \node[box] (imu) at (0,2.2) {MPU6050\\accél./gyro. (I2C)};
  \node[box] (gps) at (0,0) {GPS NEO-6M\\(UART2)};
  \node[box] (btn) at (0,-2.2) {Bouton\\d'annulation (GPIO\,33)};

  \node[esp] (mcu) at (4.8,0) {\textbf{ESP32}\\[4pt]FSM de chute\\4 phases\\[4pt]Évaluation d'état\\NORMAL / JAUNE / ROUGE};

  \node[box] (alert) at (9.6,2.2) {Buzzer\\+ LED RGB};
  \node[box] (wifi) at (9.6,0) {Wi-Fi $\rightarrow$ API REST\\$\rightarrow$ app mobile};
  \node[box] (gsm) at (9.6,-2.2) {GSM SIM800L\\$\rightarrow$ SMS};

  \draw[arrow] (imu.east) -- (mcu.west |- imu.east);
  \draw[arrow] (gps.east) -- (mcu.west |- gps.east);
  \draw[arrow] (btn.east) -- (mcu.west |- btn.east);
  \draw[arrow] (mcu.east |- alert.west) -- (alert.west);
  \draw[arrow] (mcu.east |- wifi.west) -- (wifi.west);
  \draw[arrow] (mcu.east |- gsm.west) -- (gsm.west);
\end{tikzpicture}
\caption{Architecture fonctionnelle du système SentinelBand (antenne PCB interne de l'ESP32 ; aucune antenne externe).}
\label{fig:architecture}
\end{figure}

\subsection{Justification des composants et nomenclature}

\begin{table}[H]
\centering
\caption{Nomenclature des composants du prototype SentinelBand.}
\label{tab:nomenclature}
\footnotesize
\begin{tabularx}{\textwidth}{>{\raggedright\arraybackslash}p{2.6cm} >{\raggedright\arraybackslash}X >{\raggedright\arraybackslash}X >{\raggedright\arraybackslash}p{2.6cm} c}
\toprule
\textbf{Composant} & \textbf{Fonction} & \textbf{Interface} & \textbf{Alimentation} & \textbf{Qté} \\
\midrule
ESP32 DevKit v1 & MCU principal, Wi-Fi 802.11 b/g/n, 240\,MHz dual-core & USB, GPIO, I2C, UART & 5\,V (USB ou boost) & 1 \\
\addlinespace
MPU6050 & Accéléromètre 3 axes ($\pm$8\,g, 4096\,LSB/g) et gyroscope 3 axes ($\pm$250\,°/s, 131\,LSB/(°/s)) & I2C (SDA=21, SCL=22), adresse 0x68 & 3,3\,V & 1 \\
\addlinespace
GPS NEO-6M & Géolocalisation, trames NMEA (TinyGPSPlus) & UART2 (RX=16, TX=17), 9600\,bds & 3,3\,V & 1 \\
\addlinespace
SIM800L & GSM (2G), envoi de SMS (AT+CMGS) & UART1 (RX=18, TX=19), 9600\,bds & 3,7--4,2\,V direct sur LiPo, condensateur 1000\,µF (pics de courant) & 1 \\
\addlinespace
Buzzer piézo actif & Alerte sonore JAUNE/ROUGE & GPIO\,4 & 3,3\,V & 1 \\
\addlinespace
LED RGB (cathode commune) & Indicateur visuel d'état & GPIO\,25 (R), 26 (G), 27 (B) & 3,3\,V & 1 \\
\addlinespace
Bouton poussoir NO & Annulation d'alerte & GPIO\,33 (\texttt{INPUT\_PULLUP}) & --- & 1 \\
\addlinespace
Batterie LiPo 3,7\,V & Alimentation autonome & TP4056 (charge) + MT3608 (boost 5\,V) & 3000--4000\,mAh & 1 \\
\bottomrule
\end{tabularx}
\end{table}

\textbf{Justification des choix~:}
\begin{itemize}
    \item \textbf{ESP32}~\cite{ref_esp32_datasheet} plutôt qu'Arduino Uno ou STM32~: Wi-Fi et Bluetooth intégrés, deux cœurs à 240\,MHz, mode deep sleep à faible consommation et coût réduit, ce qui en fait un choix adapté à un prototype IoT wearable. La bibliothèque HTTPClient native simplifie la communication REST.

    \item \textbf{MPU6050}~\cite{ref_mpu6050_datasheet} plutôt que l'ADXL345 (accéléromètre seul)~: capteur 6 axes largement documenté, à faible coût, en I2C standard. La plage $\pm$8\,g offre une bonne résolution (4096\,LSB/g) mais peut saturer lors d'un impact au poignet (voir section~\ref{sec:limites}).

    \item \textbf{SIM800L} pour la redondance de communication~: il permet l'envoi de SMS d'urgence hors couverture Wi-Fi. Il est toutefois limité au réseau 2G, dont la disponibilité diminue selon les pays.

    \item \textbf{GPS NEO-6M} plutôt que NEO-M8N~: compromis coût/performance suffisant pour un prototype. La bibliothèque TinyGPSPlus~\cite{ref_tinygps} décode les trames NMEA sans dépendance lourde.
\end{itemize}

% =====================================================================
%          III. MÉTHODES EXPÉRIMENTALES
% =====================================================================
\section{Méthodes expérimentales}

Le prototype utilise l'antenne PCB interne de l'ESP32 pour le Wi-Fi (2,4\,GHz, bande ISM). Les tests fonctionnels portent sur la chaîne complète~: capteurs, décision embarquée, actionneurs et transmission.

\subsection{Vérification des capteurs}

\begin{itemize}
    \item \textbf{MPU6050~:} l'initialisation écrit le registre 0x6B (Power Management) et configure les plages $\pm$8\,g (registre 0x1C, valeur 0x10) et $\pm$250\,°/s (registre 0x1B, valeur 0x00). Les données brutes sont lues par blocs de 14 octets à partir du registre 0x3B, avec les facteurs de conversion 4096\,LSB/g et 131\,LSB/(°/s). Au repos, $|\vec{a}|$ converge vers $\sim$1\,g. L'échantillonnage est fixé à 50\,Hz (20\,ms).

    \item \textbf{GPS NEO-6M~:} le tampon UART est porté à 1024 octets (\texttt{gpsSerial.\allowbreak setRxBufferSize\allowbreak(1024)}) pour éviter les pertes de trames pendant les appels HTTP. La validité du fix est testée par \texttt{gps.location.isValid()}. La commande série \texttt{g} affiche satellites, HDOP et trames valides/corrompues. En intérieur, aucune coordonnée de substitution n'est envoyée (\texttt{TEST\_INDOOR\_GPS\_COORDS~=~false}).

    \item \textbf{Module SIM800L~:} la commande série \texttt{p} exécute 6 commandes AT (\texttt{AT}, \texttt{AT+CCID}, \texttt{AT+CPIN?}, \texttt{AT+CSQ}, \texttt{AT+CREG?}, \texttt{AT+COPS?}) pour contrôler la communication, la SIM, le signal et l'enregistrement réseau. Un SMS de test est déclenché par la commande \texttt{m}~: l'envoi a été réalisé et le SMS a bien été reçu.

    \item \textbf{Température et gaz~:} aucun capteur physique. Les seuils associés ($T_{\text{jaune}} = 45$\,°C, $T_{\text{rouge}} = 60$\,°C, gaz$_{\text{jaune}} = 300$\,mV, gaz$_{\text{rouge}} = 800$\,mV) ne sont donc pas testés.
\end{itemize}

\subsection{Scénarios de test}

\begin{table}[H]
\centering
\caption{Scénarios de test fonctionnel du prototype SentinelBand.}
\label{tab:scenarios}
\footnotesize
\begin{tabularx}{\textwidth}{>{\raggedright\arraybackslash}X >{\raggedright\arraybackslash}X >{\raggedright\arraybackslash}X >{\raggedright\arraybackslash}X >{\raggedright\arraybackslash}X}
\toprule
\textbf{Scénario} & \textbf{Stimulus} & \textbf{Seuils} & \textbf{Sortie attendue} & \textbf{Statut} \\
\midrule
État normal & Mouvement modéré, $|\vec{a}| \approx 1$\,g & --- & LED verte, pas de buzzer & Testé \\
\addlinespace
Immobilité prolongée & Aucun mouvement $>$20\,s & $\Delta t > 20$\,s & JAUNE, buzzer intermittent & Testé (logs série) \\
\addlinespace
Chute simulée (série) & Commande \texttt{f} & Confiance $= 95\,\%$ injectée & JAUNE puis ROUGE, SMS, push & Testé (chaîne d'alerte) \\
\addlinespace
Chute physique réelle & Lâcher/chute contrôlée sur matelas & FSM 4 phases, conf. $\geq 60$ & JAUNE ou ROUGE selon confiance & À valider (protocole défini) \\
\addlinespace
Forçage JAUNE / ROUGE (série) & Commandes \texttt{j} / \texttt{r} & --- & Passage à l'état demandé (+ SMS en ROUGE) & Implémenté \\
\addlinespace
Annulation & Appui GPIO\,33 & Holdoff 15\,s & Retour NORMAL & Implémenté \\
\addlinespace
Escalade JAUNE$\rightarrow$ROUGE & Pas d'appui pendant 15\,s & \texttt{YELLOW\_TIMEOUT\_MS} & ROUGE + SMS & Testé (via chute simulée) \\
\addlinespace
Transmission API & HTTP POST & Code 2xx & Données en base & Testé (HTTP 200 dans les logs) \\
\addlinespace
Reconnexion Wi-Fi & Coupure du point d'accès & Tentative toutes les 8\,s & Reconnexion automatique & Testé \\
\addlinespace
Envoi SMS & Alerte ROUGE & \texttt{AT+CMGS} & SMS reçu & Testé (SMS reçu) \\
\bottomrule
\end{tabularx}
\end{table}

\subsection{Protocole de test de chute}

Pour les chutes réelles, le boîtier est fixé à la ceinture par un bandeau élastique et les chutes sont réalisées sur matelas. Trois types de chutes sont prévus~: avant (perte d'équilibre, réception sur genoux/mains), latérale (impact sur la hanche) et arrière (chute sur les fesses/dos). La commande série \texttt{s} permet de lire l'état complet (phase FSM, accélération, gyroscope), et les logs série (115200\,bauds, préfixe \texttt{[T+Xs]}) constituent la trace horodatée des essais. Ce protocole est défini~; ses résultats ne sont pas encore disponibles et le tableau~\ref{tab:chutes_reelles} est à renseigner.



% =====================================================================
%              IV. RÉSULTATS
% =====================================================================
\section{Résultats}

\subsection{Algorithme de détection de chute}

L'algorithme FSM opère en quatre phases séquentielles. Les paramètres suivants sont extraits du code embarqué (\texttt{SentinelBand\_NoSIM.ino})~:

\begin{table}[H]
\centering
\caption{Paramètres de l'algorithme FSM de détection de chute.}
\label{tab:fsm_params}
\begin{tabular}{l l c}
\toprule
\textbf{Phase} & \textbf{Paramètre} & \textbf{Valeur} \\
\midrule
\multirow{2}{*}{Phase 1 --- Pré-chute} & Seuil gyroscope & $>80$\,°/s \\
 & Fenêtre d'observation & 500\,ms \\
\addlinespace
\multirow{3}{*}{Phase 2 --- Apesanteur} & Seuil $|\vec{a}|$ max & $< 0{,}55$\,g \\
 & Durée minimale & 80\,ms \\
 & Durée maximale & 1500\,ms \\
\addlinespace
\multirow{2}{*}{Phase 3 --- Impact} & Seuil pic d'accélération & $> 2{,}3$\,g \\
 & Délai max après apesanteur & 800\,ms \\
\addlinespace
\multirow{3}{*}{Phase 4 --- Posture} & Seuil angle de gîte & $> 50$\,° pendant 5\,s \\
 & Seuil immobilité gyroscope & $< 20$\,°/s \\
 & Bande de tolérance $|\vec{a}|$ & $\bigl||\vec{a}| - 1\bigr| < 0{,}15$\,g \\
\addlinespace
Score & Confiance minimale & $\geq 60$ \\
\addlinespace
\multirow{2}{*}{Sécurité temporelle} & Timeout JAUNE$\rightarrow$ROUGE & 15\,s \\
 & Seuil immobilité & 20\,s \\
\bottomrule
\end{tabular}
\end{table}

La phase 1 (pré-chute) est obligatoire~: sans agitation gyroscopique supérieure à 80\,°/s dans la fenêtre d'observation de 500\,ms, la FSM ne progresse pas vers la phase d'apesanteur. Elle n'entre toutefois pas dans le score de confiance, qui pondère les trois autres phases~:
\begin{equation}
    C_{\text{final}} = 0{,}35 \times S_{\text{FF}} + 0{,}45 \times S_{\text{impact}} + 0{,}20 \times S_{\text{posture}}
    \label{eq:confidence}
\end{equation}
où~:
\begin{itemize}
    \item $S_{\text{FF}}$ est la moyenne de deux sous-scores. Score de durée~: 60 si 80--199\,ms, 100 si 200--599\,ms, 70 si 600--1499\,ms. Score de magnitude minimale~: 100 si $|\vec{a}|_{\min}<0{,}2$\,g, 75 si $<0{,}4$\,g, 50 sinon~;
    \item $S_{\text{impact}}$ est le score du pic à l'impact~: 30 si pic $<2{,}3$\,g, 80 si $<4$\,g, 100 si $<7$\,g, 60 au-delà. La baisse à 60 pour les pics $\geq 7$\,g peut s'interpréter comme une pénalisation de pics extrêmes, possiblement dus à un choc contre un objet plutôt qu'à une chute~; ce choix est empirique et la saturation du capteur à 8\,g limite de toute façon la mesure~;
    \item $S_{\text{posture}} = \min\!\left(\frac{\theta - 50}{40} \times 100,\ 100\right)$, où $\theta$ est l'angle de gîte en degrés.
\end{itemize}

Plusieurs mécanismes limitent les faux positifs~: reset si la durée d'apesanteur sort de l'intervalle [80~; 1500[\,ms, reset si l'impact ne survient pas dans les 800\,ms suivant l'apesanteur, et reset si la personne se relève ($\theta<30$° avec mouvement pendant $>10$\,s).



\subsection{Essais fonctionnels}

\textbf{Chute simulée (commande \texttt{f})~:} la simulation injecte une chute parfaite avec \texttt{fallConfidence}$=95\,\%$ et déclenche la séquence JAUNE$\rightarrow$ROUGE. Ce test, reproduit à plusieurs reprises, valide la chaîne d'alerte~: buzzer, LED, POST HTTP vers l'API et notification push. Il ne valide \emph{pas} la détection elle-même, puisque la FSM est court-circuitée.

\textbf{Immobilité prolongée~:} après 20\,s d'immobilité ($\texttt{gyroDps}<20$\,°/s et $\bigl||\vec{a}|-1\bigr|<0{,}15$\,g), le système passe en JAUNE avec le motif \texttt{IMMOBILITE}, observé dans les logs série.

\textbf{Reconnexion Wi-Fi~:} la déconnexion volontaire du point d'accès a confirmé une nouvelle tentative toutes les 8\,s (log \texttt{[T+Xs] Wi-Fi non connecté : nouvelle tentative}).

\textbf{Transmission télémétrique~:} le JSON transmis contient \texttt{state}, \texttt{reason}, \texttt{temp\_c}, \texttt{gas\_mv}, \texttt{acc\_g}, \texttt{gyro\_dps}, \texttt{gps\_fix}, \texttt{uptime\_s}, \texttt{fall\_phase}, \texttt{fall\_confidence}, \texttt{posture\_angle\_deg} et \texttt{source} (\texttt{"REEL"}). La réponse HTTP~200 est confirmée dans les logs (\texttt{[API] Envoi OK HTTP 200}).



\subsection{Chronologie et latences}

Les paramètres de la FSM permettent de calculer la chronologie suivante à partir du début de la chute libre (valeurs théoriques, non chronométrées)~:
\begin{itemize}
    \item \textbf{Apesanteur~:} 80 à 1500\,ms.
    \item \textbf{Impact~:} au plus 800\,ms après la fin de l'apesanteur, soit au plus 2,3\,s après le début de la chute.
    \item \textbf{Validation de posture~:} 5\,s après l'impact. L'état JAUNE est donc atteint environ \textbf{5 à 7,5\,s} après le début de la chute.
    \item \textbf{Escalade~:} 15\,s sans appui sur le bouton, soit ROUGE environ \textbf{20 à 22,5\,s} après le début de la chute.
    \item \textbf{Notification distante~:} période de télémétrie de 2\,s en alerte, traitement serveur $<$100\,ms, puis délai Expo Push Service (typiquement quelques secondes selon la documentation Expo), soit environ 2 à 7\,s après le passage en ROUGE.
\end{itemize}
Le délai total entre le début de la chute et la notification sur l'application est donc estimé à \textbf{22 à 30\,s}. Les timeouts HTTP (2500\,ms de connexion et 3000\,ms de réponse, soit 5,5\,s maximum par cycle) peuvent le rallonger. La télémétrie est émise toutes les 4\,s en NORMAL et toutes les 2\,s en JAUNE/ROUGE. Le SMS dépend de la qualité du signal GSM (typiquement 2 à 10\,s) et son succès est confirmé par la réponse \texttt{+CMGS: <ref>}.

\subsection{Architecture logicielle backend et mobile}

Le backend est construit sur Node.js ($\geq$18) avec Express~4.21.0 et PostgreSQL~16 déployé par Docker Compose. Le schéma comprend six tables~: \texttt{users}, \texttt{devices}, \texttt{colleague\_links}, \texttt{telemetry}, \texttt{device\_status} et \texttt{alerts}. L'authentification repose sur JWT (30~jours, jsonwebtoken~9.0.2)~; les clés API des boîtiers sont stockées sous forme de hachés SHA-256 et les mots de passe hachés avec bcryptjs (10 rounds). La table \texttt{device\_status} contient le dernier état connu~; \texttt{telemetry} archive l'historique (index composite sur \texttt{device\_id} et \texttt{ts DESC}).

L'application mobile (React Native~0.86.3 / Expo~57.0.27) comprend cinq écrans~: authentification, liste des collègues avec télémétrie (polling toutes les 3\,s), détail télémétrique (historique, carte GPS, alertes), demandes de liaison avec consentement explicite, et profil opérateur (état du boîtier, régénération de clé API, URL serveur). Les notifications push passent par Expo Push Service (expo-notifications~57.0.22), le JWT étant stocké via expo-secure-store~57.0.4. L'API envoie les notifications aux collègues ayant accepté le lien de surveillance et possédant un \texttt{push\_token}, uniquement lors du premier passage à l'état ROUGE.



% =====================================================================
%              V. DISCUSSION
% =====================================================================
\section{Discussion}

\subsection{Discussion du prototype wearable}

\textbf{Fiabilité des capteurs~:} le MPU6050 fournit des mesures cohérentes à 50\,Hz en I2C stable. En revanche, l'absence de capteurs de température et de gaz réels empêche toute validation de la détection d'anomalies environnementales.

\textbf{Pertinence des seuils~:} les seuils de la FSM (tableau~\ref{tab:fsm_params}) ont été fixés empiriquement, en s'inspirant de l'ordre de grandeur des valeurs rapportées dans la littérature~\cite{ref_fall_detection_threshold,ref_kangas}. La pondération plus forte de l'impact (0,45) que de l'apesanteur (0,35) traduit l'importance du pic d'accélération. Le caractère obligatoire de la phase de pré-chute (agitation gyroscopique $>80$\,°/s) filtre efficacement les faux positifs, mais peut provoquer des faux négatifs lorsqu'une chute n'implique pas de rotation marquée du la ceinture. Faute d'essais réels en nombre suffisant, les taux de faux positifs et de faux négatifs ne peuvent pas être quantifiés.

\textbf{Résolution temporelle et plage de mesure~:} à 50\,Hz, la durée minimale d'apesanteur de 80\,ms ne correspond qu'à 4 échantillons, et un impact très bref (quelques ms) peut être manqué. Un échantillonnage de 100 à 200\,Hz serait préférable. Par ailleurs, la plage $\pm$8\,g peut saturer lors d'un impact à la ceinture, ce qui tronque le pic mesuré et fausse $S_{\text{impact}}$.

\textbf{Port à la ceinture~:} le critère de posture (angle de gîte $>50$°) est naturellement conçu pour un capteur placé sur le tronc. A la ceinture, l'orientation dépend de la position du bras et non de celle du corps, ce qui peut produire à la fois des faux négatifs (bras relevé après la chute) et des faux positifs (bras en position inclinée au repos). C'est une faiblesse structurelle de l'algorithme actuel.

\textbf{Mécanismes anti-faux-positifs~:} (1)~la fenêtre d'apesanteur [80~; 1500[\,ms écarte les micro-perturbations et les chutes libres irréalistes~; (2)~le délai maximal d'impact de 800\,ms évite la confusion avec des rebonds~; (3)~le reset si la personne se relève évite les alarmes persistantes~; (4)~le délai de refroidissement après annulation (\texttt{CANCEL\_HOLDOFF\_MS}~$=15$\,s) empêche un re-déclenchement immédiat. Le critère d'immobilité de 20\,s reste cependant sensible aux faux positifs~: un travailleur qui lit un plan ou effectue une tâche statique déclenchera une alerte JAUNE.

\textbf{Latence~:} l'alerte locale (JAUNE) intervient environ 5 à 7,5\,s après le début de la chute, car la validation de posture impose 5\,s après l'impact. La notification distante n'arrive qu'après l'escalade de 15\,s, soit 22 à 30\,s au total (section~IV.C). Ce compromis privilégie la réduction des fausses alertes distantes grâce à la fenêtre d'annulation, mais retarde les secours en cas de chute réelle avec perte de conscience. Il devrait être discuté avec un responsable sécurité.

\textbf{Ergonomie~:} la carte ESP32 DevKit ($\sim$55~$\times$~28\,mm) n'est pas compatible avec un port prolongé. Un PCB sur mesure réduirait nettement l'encombrement.

\textbf{Consommation et autonomie~:} l'autonomie n'a pas été mesurée. D'après les fiches techniques, l'ESP32 en Wi-Fi actif consomme $\sim$160--260\,mA, le GPS $\sim$45\,mA et le SIM800L $\sim$7\,mA en veille (pics de 200--400\,mA en communication). Sans deep sleep, la consommation totale est estimée à 250--300\,mA, soit, avec une LiPo de 3000 à 4000\,mAh et un rendement de boost d'environ 85\,\%, une autonomie de l'ordre de 8 à 14\,h. Ces valeurs proviennent des fiches techniques et non de mesures.

\textbf{Sécurité et confidentialité~:} les clés API sont générées par \texttt{crypto.randomBytes(24)} (192 bits) et stockées hachées en SHA-256. Le lien entre collègues exige un consentement explicite (\texttt{pending} $\rightarrow$ \texttt{accepted}), en cohérence avec les principes de consentement et de minimisation du RGPD. En revanche, la communication ESP32 $\rightarrow$ API se fait en HTTP non chiffré~: la clé API transite en clair et un attaquant sur le même réseau pourrait la capturer pour usurper le boîtier.

\textbf{Perte de connexion~:} la reconnexion Wi-Fi est automatique (toutes les 8\,s) et le SMS assure une redondance. Les données non transmises pendant une déconnexion sont perdues (pas de file d'attente locale).

\subsection{Limites}
\label{sec:limites}

\begin{enumerate}
    \item \textbf{Capteurs de température et de gaz non implémentés~:} valeurs constantes, branches de détection non validées (LM35 et MQ-2 à ajouter).
    \item \textbf{Validation statistique absente~:} seuils empiriques, aucune matrice de confusion. Une campagne structurée (au moins 30 chutes par type, plusieurs porteurs, activités quotidiennes) est nécessaire pour estimer sensibilité et spécificité.
    \item \textbf{Algorithme validé uniquement par injection~:} la chaîne d'alerte est testée, pas la détection sur chutes réelles.
    \item \textbf{Échantillonnage et plage de mesure~:} 50\,Hz (4 échantillons dans 80\,ms) et saturation possible à $\pm$8\,g.
    \item \textbf{Critère de posture inadapté à la ceinture~:} voir discussion.
    \item \textbf{Autonomie estimée, non mesurée.}
    \item \textbf{Format non wearable~:} carte de développement, preuve de concept de laboratoire.
    \item \textbf{Communication non chiffrée et sans tampon local~:} HTTP en clair, pertes de données hors connexion.
    \item \textbf{SIM800L limité à la 2G}, dont la couverture disparaît dans plusieurs pays.
    \item \textbf{Environnement de test non contrôlé~:} essais en bureau/laboratoire, sans étude des interférences 2,4\,GHz, de la couverture GPS en intérieur ni du GSM d'un site industriel réel.
\end{enumerate}

% =====================================================================
%              VI. CONCLUSION
% =====================================================================
\section{Conclusion}

Ce travail a porté sur la conception et l'évaluation préliminaire d'un ceinture connecté de surveillance de la sécurité des travailleurs industriels. Le prototype SentinelBand repose sur un ESP32, un MPU6050 ($\pm$8\,g, 50\,Hz), un GPS NEO-6M, un module GSM SIM800L et des actionneurs d'alerte (buzzer, LED RGB). Une FSM de détection de chute en quatre phases avec score de confiance pondéré (équation~\ref{eq:confidence}, seuil de décision à 60) a été implémentée et instrumentée. Une chaîne de supervision complète a été développée~: API REST Node.js/PostgreSQL, application React Native avec notifications push et liaison entre collègues avec consentement explicite.

Les essais par injection de chute ont validé la chaîne d'alerte (JAUNE, ROUGE, transmission à l'API, notification push) ainsi que la reconnexion Wi-Fi automatique. L'algorithme de détection n'a en revanche pas encore été validé sur des chutes réelles, et ses performances (sensibilité, spécificité) restent inconnues. Les autres limites majeures sont l'absence de capteurs de température et de gaz, l'échantillonnage à 50\,Hz, l'inadéquation du critère de posture au port à la ceinture, l'absence de mesure d'autonomie et la transmission non chiffrée.

Deux perspectives prioritaires se dégagent~: (1)~une campagne de validation sur chutes réelles et activités quotidiennes, avec relevé des traces IMU et calcul de sensibilité et de spécificité, assortie d'un échantillonnage plus élevé et d'une adaptation du critère de posture~; (2)~la miniaturisation sur PCB sur mesure, l'optimisation de la consommation (deep sleep, télémétrie adaptative), l'ajout des capteurs LM35 et MQ-2 et la sécurisation des échanges par HTTPS (mbedTLS sur ESP32). L'intégration d'une antenne textile externe, non réalisée ici, constitue un travail futur de conception radiofréquence.

% =====================================================================
%              RÉFÉRENCES
% =====================================================================
\newpage
\begin{thebibliography}{99}

\bibitem{ref_wearable_survey}
F.~J.~Dian, R.~Vahidnia, and A.~Rahmati,
``Wearables and the Internet of Things (IoT), Applications, Opportunities, and Challenges: A Survey,''
\textit{IEEE Access}, vol.~8, pp.~69200--69211, 2020.
DOI~: \url{https://doi.org/10.1109/ACCESS.2020.2986329}

\bibitem{ref_lone_worker}
European Agency for Safety and Health at Work (EU-OSHA),
``Lone Workers --- Health and Safety Guidance,''
Publication Office of the European Union, 2021.
\url{https://osha.europa.eu/en/tools-and-resources/publications/lone-workers}

\bibitem{ref_bayar2023}
D.~Bayar, C.~Aridici, Y.~S.~Metin, O.~Za, and R.~Edizkan,
``Elderly Fall Detection System with ESP32 Module and Edge Impulse Studio,''
in \textit{Proc. 7th Int. Symp. Innovative Approaches in Smart Technologies (ISAS)}, Istanbul, 2023.
DOI~: \url{https://doi.org/10.1109/isas60782.2023.10391675}

\bibitem{ref_fall_detection_threshold}
N.~Noury, T.~Hervé, V.~Rialle, G.~Virone, E.~Mercier, G.~Morey, A.~Moro, and T.~Porcheron,
``Monitoring behavior in home using a smart fall sensor and position sensors,''
in \textit{Proc. 1st Annual Int. IEEE-EMBS Special Topic Conf. Microtechnologies in Medicine \& Biology}, 2000, pp.~607--610.
DOI~: \url{https://doi.org/10.1109/MMB.2000.893856}

\bibitem{ref_kangas}
M.~Kangas, A.~Konttila, P.~Lindgren, I.~Winblad, and T.~Jämsä,
``Comparison of low-complexity fall detection algorithms for body attached accelerometers,''
\textit{Gait \& Posture}, vol.~28, no.~2, pp.~285--291, 2008.
DOI~: \url{https://doi.org/10.1016/j.gaitpost.2008.01.003}

\bibitem{ref_mauldin}
T.~Mauldin, M.~Canby, V.~Metsis, A.~Ngu, and C.~Rivera,
``SmartFall: A Smartwatch-Based Fall Detection System Using Deep Learning,''
\textit{Sensors}, vol.~18, no.~10, p.~3363, 2018.
DOI~: \url{https://doi.org/10.3390/s18103363}

\bibitem{ref_esp32_datasheet}
Espressif Systems,
``ESP32 Series Datasheet,''
\url{https://www.espressif.com/sites/default/files/documentation/esp32_datasheet_en.pdf}, consulté le 21~octobre~2026.

\bibitem{ref_mpu6050_datasheet}
InvenSense (TDK),
``MPU-6000 and MPU-6050 Product Specification, Revision 3.4,''
Document PS-MPU-6000A-00, 2013.
\url{https://invensense.tdk.com/products/motion-tracking/6-axis/mpu-6050/}

\bibitem{ref_tinygps}
M.~Hart,
``TinyGPS++: A GPS library for Arduino,''
\url{https://github.com/mikalhart/TinyGPSPlus}

\end{thebibliography}

% =====================================================================
%              ANNEXES
% =====================================================================
\newpage
\appendix

\section{Annexe A --- Prototype embarqué}
\label{ann:prototype}

\subsection*{A.1 --- Schéma de brochage ESP32}

\begin{table}[H]
\centering
\caption{Brochage du prototype SentinelBand sur ESP32.}
\small
\begin{tabular}{l l l l}
\toprule
\textbf{GPIO ESP32} & \textbf{Composant} & \textbf{Signal} & \textbf{Note} \\
\midrule
21 & MPU6050 & SDA (I2C) & Pull-up 4,7\,k$\Omega$ \\
22 & MPU6050 & SCL (I2C) & Pull-up 4,7\,k$\Omega$ \\
16 & GPS NEO-6M & RX UART2 $\leftarrow$ TX GPS & Tampon 1024 octets \\
17 & GPS NEO-6M & TX UART2 $\rightarrow$ RX GPS & --- \\
18 & SIM800L & RX UART1 $\leftarrow$ TX SIM & TX SIM800L $\approx$2,8\,V, lisible en 3,3\,V \\
19 & SIM800L & TX UART1 $\rightarrow$ RX SIM & Diviseur conseillé sur cette ligne \\
4 & Buzzer piézo actif & Sortie numérique & LOW = silencieux \\
25 & LED RGB & Rouge & Cathode commune, 100\,$\Omega$ \\
26 & LED RGB & Vert & Cathode commune, 100\,$\Omega$ \\
27 & LED RGB & Bleu & Cathode commune, 100\,$\Omega$ \\
33 & Bouton poussoir NO & \texttt{INPUT\_PULLUP} & Actif à l'état bas, debounce 300\,ms \\
GND & MPU6050 AD0 & GND & Adresse I2C 0x68 \\
\bottomrule
\end{tabular}
\end{table}

\subsection*{A.2 --- Organigramme du firmware}

La boucle principale (\texttt{loop()}) est non bloquante et appelle successivement~:
\begin{enumerate}
    \item \texttt{readGps()} --- lecture UART, décodage NMEA par TinyGPSPlus
    \item \texttt{readSensors()} --- lecture MPU6050 à 50\,Hz, suivi d'immobilité, appel de \texttt{runFallFSM()}
    \item \texttt{checkButton()} --- annulation d'alerte sur appui du GPIO\,33 (debounce 300\,ms)
    \item \texttt{checkSerial()} --- commandes de débogage
    \item \texttt{evaluateState()} --- transitions NORMAL/JAUNE/ROUGE, SMS au passage en ROUGE
    \item \texttt{updateBuzzer()} et \texttt{updateLed()} --- actionneurs selon l'état
    \item \texttt{telemetryTick()} --- POST HTTP si la période est écoulée (4\,s normal, 2\,s alerte) et Wi-Fi connecté
    \item Log série à 1\,Hz via \texttt{printStatus()}
\end{enumerate}

\subsection*{A.3 --- Commandes série de débogage (115200\,bauds)}

\begin{table}[H]
\centering
\caption{Commandes disponibles sur le port USB de l'ESP32.}
\small
\begin{tabular}{l l}
\toprule
\textbf{Touche} & \textbf{Action} \\
\midrule
\texttt{c} ou \texttt{n} & Annuler l'alerte en cours, retour à NORMAL \\
\texttt{f} & Simuler une chute (confiance 95\,\%) \\
\texttt{j} & Forcer l'état JAUNE (\texttt{TEST\_MANUEL\_JAUNE}) \\
\texttt{r} & Forcer l'état ROUGE (\texttt{TEST\_MANUEL\_ROUGE}) + SMS \\
\texttt{t} & Forcer l'envoi immédiat de télémétrie \\
\texttt{s} & Afficher le statut complet \\
\texttt{g} & Diagnostic GPS \\
\texttt{m} & Envoyer un SMS de test \\
\texttt{p} & Diagnostic SIM800L (6 commandes AT) \\
\texttt{AT...} & Envoyer une commande AT directement au SIM800L \\
\bottomrule
\end{tabular}
\end{table}

\subsection*{A.4 --- Versions et dépendances}

\begin{table}[H]
\centering
\caption{Versions des bibliothèques et outils utilisés.}
\small
\begin{tabular}{l l l}
\toprule
\textbf{Composant} & \textbf{Technologie} & \textbf{Version} \\
\midrule
Firmware ESP32 & Arduino Core for ESP32 (Espressif) & $\geq$~2.0.x \\
Bibliothèque GPS & TinyGPSPlus (Mikal Hart) & 1.0.3 \\
Communication HTTP & HTTPClient (native Arduino ESP32) & native \\
Backend API & Node.js / Express & $\geq$18 / 4.21.0 \\
Base de données & PostgreSQL & 16 \\
Driver BDD & pg (node-postgres) & 8.12.0 \\
Authentification & jsonwebtoken / bcryptjs & 9.0.2 / 2.4.3 \\
Application mobile & React Native / Expo & 0.86.3 / 57.0.27 \\
Navigation & React Navigation & 7.x \\
Notifications push & expo-notifications & 57.0.22 \\
Stockage sécurisé & expo-secure-store & 57.0.4 \\
Infrastructure BDD & Docker Compose + PostgreSQL & postgres:16 \\
\bottomrule
\end{tabular}
\end{table}

\subsection*{A.5 --- Code source}

Le code complet (firmware ESP32, backend, application mobile) est disponible sur~:
\begin{center}
\url{https://github.com/AdemAboueb7/sentinelband.git}
\end{center}


\end{document}
